\section{Elektr maydon. Elektr maydon kuchlanganligi}


\begin{enumerate}

    \item Zaryadlangan zarracha vakuumning qandaydir nuqtasida 60 V/m kuchlanganlik hosil qiladi. Agar butun sistema dielektrik singdiruvchanligi 2 bo’lgan kerosin ichiga joylashtirilsa, shu nuqtaga o’rnatilgan 5 nC zaryadga qanday kuch (nN da) ta’sir qiladi? 
    \begin{enumerate}[label=\Alph*)]
        \item 300
        \item 150 % To'g'ri
        \item 75
        \item 60
    \end{enumerate}

    \item Kuchlanganlik vektori vertikal yuqoriga yo’nalgan bir jinsli elektr maydonda massasi 0,03 mkg bo’lgan 3 pC zaryadli chang zarrachasi muvozanat holatida turibdi. Maydon kuchlanganligini aniqlang ($g=10 \, m/s^2$). 
    \begin{enumerate}[label=\Alph*)]
        \item 100 % To'g'ri
        \item 10
        \item 1000
        \item 1
    \end{enumerate}

    \item Miqdori 20 kV/m bo’lgan kuchlanganlik vektori vertikal pastga yo’nalgan bir jinsli elektr maydonda 0,1 kg massali va 0,2 mC zaryadli sharcha ipak ipga osilgan. Ipning taranglik kuchini toping ($g=10 \, m/s^2$). 
    \begin{enumerate}[label=\Alph*)]
        \item 1
        \item 3
        \item 5 % To'g'ri
        \item 10
    \end{enumerate}

    \item Agar ipga osilgan 0,1 kg massadagi 10 mkC zaryadli sharcha 200 kV/m kuchlanganlikka ega bo’lgan bir jinsli elektr maydonda joylashtirilsa, ipning taranglik kuchi necha marta ortadi? Kuchlanganlik vektori vertikal pastga yo’nalgan ($g=10 \, m/s^2$). 
    \begin{enumerate}[label=\Alph*)]
        \item 2
        \item 3 % To'g'ri
        \item 4
        \item 1,5
    \end{enumerate}

    \item Massasi 4,5 g bo’lgan 0,1 mkC zaryadli sharcha $800 \, kg/m^3$ zichlikdagi moy ichiga joylashtirilgan. Sharcha materialining zichligi $1500 \, kg/m^3$. Sharcha muvozanatda bo'lishi uchun elektr maydonning kuchlanganligi qanday (kV/m da) bo’lishi kerak? ($g=10 \, m/s^2$). 
    \begin{enumerate}[label=\Alph*)]
        \item 210 % To'g'ri
        \item 150
        \item 180
        \item 240
    \end{enumerate}

    \item Ipak ipga osilgan kichik sharcha 49 nC zaryadga ega. 100 kV/m kuchlanganlikli gorizontal elektr maydon ta'sirida ip tangensi 0,125 bo’lgan burchakka og’di. Sharchaning massasini (g da) toping ($g=10 \, m/s^2$). 
    \begin{enumerate}[label=\Alph*)]
        \item 8
        \item 2
        \item 5
        \item 4 % To'g'ri
    \end{enumerate}

    \item Kuchlanganligi 1000 V/m bo’lgan bir jinsli elektr maydon ta’siri ostida massasi 0,1 g bo’lgan 4 mkC zaryadli zarracha erishadigan tezlanishning miqdorini toping. Og’irlik kuchini hisobga olmang. 
    \begin{enumerate}[label=\Alph*)]
        \item 4
        \item 40 % To'g'ri
        \item 400
        \item 25
    \end{enumerate}

    \item Kuchlanganligi 20 kV/m bo’lgan bir jinsli elektr maydonida massasi 0,01 kg bo’lgan 1 mkC zaryadli sharcha qanday tezlanish bilan tushishini aniqlang. Kuchlanganlik vektori vertikal yuqoriga yo’nalgan. ($g=10 \, m/s^2$). 
    \begin{enumerate}[label=\Alph*)]
        \item 8 % To'g'ri
        \item 10
        \item 2
        \item 12
    \end{enumerate}

    \item Jismga 70 nC zaryad berilganda u yer sirti yaqinida tushishning 10 sekundi ichida zaryadsiz bo’lgan paytidagidan 5 cm ga ko’proq yo’lni bosib o’tdi. Agar elektr maydon kuchlanganligi 100 V/m bo’lsa, jismning massasi (g da) qanchaga teng? 
    \begin{enumerate}[label=\Alph*)]
        \item 14
        \item 7 % To'g'ri
        \item 3,5
        \item 10
    \end{enumerate}

    \item 1 mg massali chang zarrachasi havoda 0,2 m/s o’zgarmas tezlik bilan tushadi. Agar zarracha kuchlanganligi 10 kV/m bo’lgan elektr maydonda joylashtirilib, unga 1,2 nC zaryad berilsa, u qanday barqaror tezlik (cm/s da) bilan ko’tariladi? ($g=10 \, m/s^2$). 
    \begin{enumerate}[label=\Alph*)]
        \item 2
        \item 8
        \item 4 % To'g'ri
        \item 12
    \end{enumerate}

    \item 5 mg massali zaryadlanmagan chang zarrachasi havoda 15 cm/s o’zgarmas tezlik bilan tushadi. Agar zarracha kuchlanganligi 3 kV/m bo’lgan gorizontal elektr maydonda joylashtirilib, unga 40 nC zaryad berilsa, u qanday barqaror tezlik (cm/s da) bilan harakatlanadi? 
    \begin{enumerate}[label=\Alph*)]
        \item 39 % To'g'ri
        \item 25
        \item 45
        \item 15
    \end{enumerate}

    \item 100 km/s tezlik bilan harakatlanayotgan proton kuchlanganligi 50 V/m bo’lgan elektr maydonga maydon kuch chiziqlari yo’nalishiga qarama-qarshi yo’nalishda uchib kiradi. Necha mikrosekunddan so’ng protonning tezligi nolga teng bo’ladi? Proton zaryadining massasiga nisbati $10^8 \, C/kg$. 
    \begin{enumerate}[label=\Alph*)]
        \item 10
        \item 20 % To'g'ri
        \item 40
        \item 5
    \end{enumerate}

    \item Nuqtaviy zaryad vakuumning qandaydir nuqtasida 600 V/m kuchlanganlikdagi maydonni yuzaga keltiradi. Agar zaryad 5 marta ortib, uning atrofidagi bo’shliq dielektrik singdiruvchanligi 2 bo’lgan kerosin bilan to’ldirilsa, shu nuqtada maydon kuchlanganligi qanday bo’ladi? 
    \begin{enumerate}[label=\Alph*)]
        \item 3000
        \item 600
        \item 1500 % To'g'ri
        \item 1200
    \end{enumerate}

    \item Zaryad tomonidan hosil qilinayotgan maydon kuchlanganligi 10 cm masofada 800 V/m ga teng. Zaryaddan 20 cm masofada bo’lgan nuqtaning maydon kuchlanganligini toping. 
    \begin{enumerate}[label=\Alph*)]
        \item 400
        \item 200 % To'g'ri
        \item 100
        \item 800
    \end{enumerate}

    \item 4 nC miqdordagi, turli ishorali ikkita bir xil nuqtaviy zaryad bir-biridan 60 sm masofada joylashgan. Zaryadlarni tutashtiruvchi kesmaning o’rtasida joylashgan nuqtaning maydon kuchlanganligini toping. $k=9 \cdot 10^9 \, m/F$. 
    \begin{enumerate}[label=\Alph*)]
        \item 400
        \item 1600
        \item 800 % To'g'ri
        \item 200
    \end{enumerate}

    \item Ikki musbat nuqtaviy zaryad orasidagi masofa 8 cm. Zaryadlarni birlashtiruvchi to’g’ri chiziqda, birinchi zaryaddan 6 cm masofada maydon kuchlanganligi nolga teng. Birinchi zaryad miqdorining ikkinchi zaryad miqdoriga nisbatini toping. 
    \begin{enumerate}[label=\Alph*)]
        \item 3
        \item 9 % To'g'ri
        \item 4
        \item 16
    \end{enumerate}

    \item Agar 2 nC va –4 nC nuqtaviy zaryadlarni birlashtiruvchi kesmaning o’rtasida yotuvchi nuqtada faqat birinchi zaryad hosil qiladigan maydon kuchlanganligi 2 V/m ga teng bo’lsa, shu nuqtada ikkala zaryad hosil qiladigan maydon kuchlanganligining miqdorini toping. 
    \begin{enumerate}[label=\Alph*)]
        \item 6 % To'g'ri
        \item 2
        \item 4
        \item 8
    \end{enumerate}

    \item Musbatining miqdori manfiysining miqdoridan 2,25 marta katta bo’lgan ikkita turli nomdagi nuqtaviy zaryadlar bor. Zaryadlar orasidagi masofa manfiy zaryaddan maydon kuchlanganligi nolga teng bo’lgan nuqtagacha bo’lgan masofadan necha marta kichik? 
    \begin{enumerate}[label=\Alph*)]
        \item 0,5
        \item 1
        \item 3
        \item 2 % To'g'ri
    \end{enumerate}

    \item 64 nC va –48 nC nuqtaviy zaryadlar orasidagi masofa 10 cm ga teng. Birinchi zaryaddan 8 cm va ikkinchisidan 6 cm uzoqlikda bo’lgan nuqtadagi maydon kuchlanganligini (kV/m da) aniqlang. $k=9 \cdot 10^9 \, m/F$. 
    \begin{enumerate}[label=\Alph*)]
        \item 100
        \item 150 % To'g'ri
        \item 200
        \item 250
    \end{enumerate}

    \item Miqdori bir xil 36 nC, ishorasi esa turli xil bo’lgan nuqtaviy zaryadlar tomoni 2 m bo’lgan teng tomonli uchburchakning ikki uchida joylashgan. Uchburchakning uchinchi uchidagi elektr maydon kuchlanganligini aniqlang. $k=9 \cdot 10^9 \, m/F$. 
    \begin{enumerate}[label=\Alph*)]
        \item 162
        \item 81 % To'g'ri
        \item 40,5
        \item 27
    \end{enumerate}

    \item Miqdori bir xil 5 nC, ishorasi esa turli xil bo’lgan nuqtaviy zaryadlar bir-biridan 2,4 m masofada joylashgan. Har bir zaryaddan 3 m uzoqlikda bo’lgan nuqtadagi elektr maydon kuchlanganligini toping. $k=9 \cdot 10^9 \, m/F$. 
    \begin{enumerate}[label=\Alph*)]
        \item 2
        \item 8
        \item 4 % To'g'ri
        \item 10
    \end{enumerate}

\end{enumerate}